\documentclass[11pt]{article}

\usepackage[T1]{fontenc}
\usepackage[utf8]{inputenc}
\usepackage{lmodern}
\usepackage[a4paper,margin=25mm]{geometry}
\usepackage{amsmath,amssymb,amsthm,mathtools}
\usepackage{microtype}
\usepackage{xcolor}
\usepackage{enumitem}
\usepackage{booktabs}
\usepackage{tabularx}
\usepackage{longtable}
\usepackage{hyperref}
\setlength{\emergencystretch}{3em}

\hypersetup{
  colorlinks=true,
  linkcolor=blue!55!black,
  citecolor=blue!55!black,
  urlcolor=blue!55!black
}

\newtheorem{theorem}{Theorem}[section]
\newtheorem{proposition}[theorem]{Proposition}
\newtheorem{lemma}[theorem]{Lemma}
\newtheorem{corollary}[theorem]{Corollary}
\newtheorem{definition}[theorem]{Definition}
\theoremstyle{remark}
\newtheorem{remark}[theorem]{Remark}

\newcommand{\A}{\mathcal A}
\newcommand{\B}{\mathcal B}
\newcommand{\F}{\mathcal F}
\newcommand{\Hh}{\mathcal H}
\newcommand{\K}{\mathcal K}
\newcommand{\V}{\mathcal V}
\newcommand{\C}{\mathbb C}
\newcommand{\R}{\mathbb R}
\newcommand{\Z}{\mathbb Z}
\newcommand{\id}{\mathbf 1}
\newcommand{\Tr}{\operatorname{Tr}}
\newcommand{\Ad}{\operatorname{Ad}}
\newcommand{\CAR}{\operatorname{CAR}}
\newcommand{\Ind}{\operatorname{Ind}}
\newcommand{\supp}{\operatorname{supp}}
\newcommand{\HS}{\mathrm{HS}}
\newcommand{\Dfin}{\mathcal D_{\mathrm{fin}}}
\newcommand{\norm}[1]{\left\lVert#1\right\rVert}
\newcommand{\ip}[2]{\left\langle#1,#2\right\rangle}
\newcommand{\mP}{\textsf{PROVED}}
\newcommand{\mR}{\textsf{REDUCED-TO-CITED}}
\newcommand{\mG}{\textsf{GAP}}

\title{The Charged MMST Scaling-Limit Theorem:\\
A Complete In-House Proof Attempt for
\texttt{SEAM.MMST.TYPEIII.CHARGED.01}}
\author{Proof memorandum}
\date{30 August 2026}

\begin{document}
\maketitle

\begin{abstract}
We attempt the complete scaling-limit theorem for the concrete order-four
charged field
\[
 \Psi_{\lambda,N}(f)
 =\varepsilon_\lambda:\!\exp(i\lambda\cdot\phi_N)\!:(f),
 \qquad
 \lambda=(\omega_s,\omega_f),\quad
 \norm{\lambda}^2=2,\quad h_\lambda=1.
\]
The algebraic limit-exchange statement is proved in full for an arbitrary
finite group: equivariant injective inductive limits commute isometrically
with crossed products, coefficient and fixed-point expectations commute
with the limit, and their Pimsner--Popa constants remain uniform.  Assuming
the charged field and local algebras exist, the identification of the
extension with the holomorphic $(E_8)_1$ net is reduced, with every
hypothesis displayed, to the Kawahigashi--Longo--M\"uger index theorem,
the rank-eight even-unimodular lattice classification, the
Dong--Mason central-charge-eight theorem, and the
Carpi--Kawahigashi--Longo--Weiner VOA--net correspondence.

The uniform Green-form problem is closed here for the edge-projected
frozen collar, with the explicit all-volume constant
$K_G=\pi^2/4$.  The proof combines an exact particle--hole count,
the sharp fundamental-domain sine inequality, and analytic perturbation
of the isolated edge doublet.  For the edge-projected boson modes we
also construct exact compatible dyadic Fock isometries and prove the
charged-field graph defect with rate $N^{-2}$, including adjoints.
The remaining hard lattice-to-field step is microscopic local-algebra
exhaustion.  The frozen computations
verify the proposed two-edge remainder bound
$\norm{R_N}_{\HS}\le2.11$ only for
$N=16,24,32,48,64,96$ (maximum $2.105079358$), not for all $N$.
Consequently the master theorem remains conditional on one analytic
gap box, stated quantitatively at the end.  This is a full proof attempt,
not a claim that finite observations establish a type-$\mathrm{III}$
scaling limit.
\end{abstract}

\tableofcontents

\section{Verdict and logical form}

\begin{theorem}[Charged MMST master theorem: exact proved form]
\label{thm:master}
Let the finite collars, scaling maps, neutral current limit and charged
operators be as in Section~\ref{sec:setup}.  Suppose the scaling algebra
is exhausted by the neutral currents and
the four charged grades as specified in \textup{(ALG-EXH)} there.  Then:
\begin{enumerate}[label=\textup{(L\arabic*)}]
\item For every $f\in C^\infty(S^1)$ there are $p<\infty$ and
  $C_f<\infty$, independent of $N$, such that on the finite-energy core
  \[
    \norm{\Psi_{\lambda,N}(f)\xi}
    \le C_f\norm{(1+L_{0,N})^p\xi}.
  \tag{1.1}\label{eq:L1}
  \]
\item After the canonical scaling identifications,
  $\Psi_{\lambda,N}(f)$ converges strongly on the common finite-energy
  core to a nonzero closable field $\Psi_\lambda(f)$.  The generated local
  algebras converge to a local, conformal extension net
  $I\mapsto\B(I)$ of the neutral net $I\mapsto\A(I)$.
\item The finite $\Z_4$ crossed products, the canonical coefficient
  expectations, and the fixed-point expectations commute with the
  inductive/scaling limit.  The limiting Pimsner--Popa bound is still
  $E(x)\ge\frac14x$ for $x\ge0$.
\item The resulting extension is the strongly local lattice net
  $\A_{V_{E_8}}$, equivalently the level-one $E_8$ current net.  It is
  holomorphic, has $\mu=1$, central charge $8$, vacuum character
  $E_4/\eta^8$, and local algebras equal to the hyperfinite
  type-$\mathrm{III}_1$ factor.
\end{enumerate}
\end{theorem}

\begin{proof}
L1 follows from Theorem~\ref{thm:L1reduction}.  L2 follows from
Theorem~\ref{thm:L2reduction}.  L3 is the unconditional
Theorem~\ref{thm:crossedlimit}.  L4 follows from
Theorem~\ref{thm:identification}.  The Green estimate used by L1 is the
all-volume Theorem~\ref{thm:UGF}; the graph estimate used by L2 is the
all-volume dyadic Theorem~\ref{thm:TELa}.  The condition still used by L2
is precisely \textup{(ALG-EXH)} in the final gap box; it is not silently
replaced by a finite-volume fit.
\end{proof}

\begin{remark}[What is and is not proved]
The implication in Theorem~\ref{thm:master} is proved.  Uniform Green
domination and local Sobolev convergence are proved in
Theorem~\ref{thm:UGF}, and edge-mode graph convergence is proved in
Theorem~\ref{thm:TELa}.  The antecedent \textup{(ALG-EXH)} is not proved
for the frozen microscopic collar.  Therefore
this memorandum does not close
\texttt{SEAM.MMST.TYPEIII.CHARGED.01}.  It proves L3, completes the
classical reduction of L4, closes L1 and the edge-projected TEL-a part
of L2, and isolates the remaining microscopic local-algebra statement.
\end{remark}

\section{Setup: the frozen collar and charged generator}
\label{sec:setup}

\subsection{The finite CAR collar}

For circumference $N$ and fixed open width $N_y=8$, the one-particle
Qi--Wu--Zhang collar is
\[
 h(k_x,k_y)
 =\sin k_x\,\sigma_x+\sin k_y\,\sigma_y+
   (1-\cos k_x-\cos k_y)\sigma_z .
\tag{2.1}\label{eq:qwz}
\]
The $r$th deck sector, $r\in\Z_4$, has boundary phase $i^r$, hence
longitudinal momenta
\[
 p_n^{(r)}=\frac{2\pi}{N}\left(n+\frac r4\right).
\tag{2.2}\label{eq:momentum}
\]
Sixteen Majorana copies produce the carrier central charge
$16\cdot\frac12=8$.  We write $\mathfrak h_N$ for the one-particle space,
$H_N^{(r)}$ for the corresponding one-particle matrix, and
\[
 C_N^{(r)}
 =\chi_{(-\infty,0)}(H_N^{(r)})
  \frac12\chi_{\{0\}}(H_N^{(r)})
\tag{2.3}\label{eq:covariance}
\]
for the particle--hole symmetric covariance.  The finite observable
algebra is a matrix CAR algebra.  Its low-energy GNS space is denoted
$\Hh_N$, and $L_{0,N}$ denotes the shifted, nonnegative edge-energy
operator in the scaling normalization.

\begin{remark}[Necessary correction: there is no uniform full gap]
\label{rem:gapcorrection}
The topological collar has chiral edge levels with spacing $O(N^{-1})$.
Thus a claim of a spectral gap
$\inf_N\operatorname{gap}(H_N)>0$ is false.  The usable hypothesis is a
\emph{uniform bulk gap after removing the edge band}.  If $Q_N$ is the
edge-band projection, the correct statement is
\[
 \operatorname{spec}\bigl((1-Q_N)H_N(1-Q_N)\bigr)
 \cap(-\delta_{\rm bulk},\delta_{\rm bulk})=\varnothing
\tag{2.4}\label{eq:bulkgap}
\]
for one $\delta_{\rm bulk}>0$ independent of $N$.  This controls bulk
resolvents but does not control the massless longitudinal Green function
by itself.
\end{remark}

\subsection{The exact order-four vector}

In rational coordinates,
\[
 \omega_s=\left(\frac12,\frac12,\frac12,\frac12,\frac12\right),
 \qquad
 \omega_f=\left(\frac34,-\frac14,-\frac14,-\frac14\right),
 \qquad
 \lambda=(\omega_s,\omega_f).
\tag{2.5}\label{eq:lambda}
\]
The exact finite calculation gives
\[
 \norm{\omega_s}^2=\frac54,\quad
 \norm{\omega_f}^2=\frac34,\quad
 \norm{\lambda}^2=2,\quad h_\lambda=1,
\tag{2.6}\label{eq:lambda-norm}
\]
and
\[
 [\lambda]^4=[0],\qquad [\lambda]^2=[v].
\tag{2.7}\label{eq:lambda-order}
\]
The norm-two vectors in
$D_5\oplus A_3+k\lambda$, $k=0,1,2,3$, have census
\[
 [52,64,60,64],
\tag{2.8}\label{eq:census}
\]
so the four cosets contain $240$ roots and the two odd grades contain the
full $128$-dimensional spinor contribution.

\subsection{The finite charged field}

Let $\phi_N=(\phi_N^1,\ldots,\phi_N^8)$ be the bosonic current obtained
from the sixteen Majorana copies by edge-projected collar bosonization,
with zero mode fixed by the $\Z_4$ charge sector.  Each complex-fermion
pair gives one boson.  The longitudinal derivative uses the frozen
lattice symbol
\[
 \omega_N(m)=\frac{N}{\pi}\left|\sin\frac{\pi m}{N}\right|,
 \qquad 0<|m|\le\frac N2.
\tag{2.9}\label{eq:latticesymbol}
\]
The edge projection is the analytic top (or bottom) eigenbranch through
the zero-energy crossing; massive bulk-complement modes are not included
in $\phi_N$.  For a fixed normalized
two-cocycle $\varepsilon_\lambda$, define on the finite-energy vectors
\[
 \Psi_{\lambda,N}(f)
 =\sum_{m}\widehat f_m\,
   \varepsilon_\lambda
   :\!\exp(i\lambda\cdot\phi_N)\!:_m .
 \tag{2.10}\label{eq:chargedfield}
\]
Normal ordering is with respect to $C_N^{(0)}$.  Formula
\eqref{eq:chargedfield} is meaningful at every finite $N$.  The theorem
does not follow merely from that fact: the normal-ordering constant and
the graph norm of the exponential can grow with $N$.

\begin{definition}[Scaling convergence used below]
Let $V_N:\Hh_N\to\Hh_\infty$ be the canonical isometries of the chosen
operator-algebraic renormalization system.  A sequence of closable
operators $X_N$ converges on the finite-energy core
\[
 \Dfin=\bigcup_{E<\infty}\id_{[0,E]}(L_0)\Hh_\infty
 \tag{2.11}\label{eq:core}
\]
if for every $\xi\in\Dfin$, represented at sufficiently large scale,
\[
 V_NX_NV_N^*\xi\longrightarrow X\xi,\qquad
 V_NX_N^*V_N^*\xi\longrightarrow X^*\xi.
 \tag{2.12}\label{eq:graphconv}
\]
The adjoint clause is included because bare pointwise convergence of
unbounded fields is too weak to determine local von Neumann algebras.
\end{definition}

\section{Classical inputs, quoted in the forms used}
\label{sec:classical}

\begin{theorem}[Osborne--Stottmeister free-fermion scaling theorem,
form used here]\label{thm:OS}
For the massless chiral lattice fermion systems constructed by
operator-algebraic renormalization, with the regular scaling functions
and compatible ground states specified in their construction, the
renormalized Koo--Saleur approximants converge strongly on the
finite-particle core to the corresponding Virasoro generators.
The smeared approximants converge on that core, and the same argument
applies to level-one Wess--Zumino--Witten currents which are fermion
bilinears.  Their dynamical correlation functions converge uniformly on
compact time intervals under the stated regularity assumptions on the
smearing functions.
\end{theorem}

This is the content used from Osborne--Stottmeister,
\emph{Conformal field theory from lattice fermions},
Theorems 4.11, 4.16, 5.1, 6.1 and 6.3
\cite{OsborneStottmeister}.  It supplies the neutral bilinear current net
and the common stress tensor.  It does not contain
\eqref{eq:chargedfield}: a spin/lattice exponential changes sectors and
is not a fermion bilinear.

\begin{theorem}[Energy bounds and smearing for unitary VOAs]
\label{thm:CKLWenergy}
Let $V$ be a unitary VOA with conformal Hamiltonian $L_0$, and let
$a\in V$ be homogeneous.  The field $Y(a,z)=\sum_na_nz^{-n-d_a}$
is polynomially energy bounded if there exist $M>0$ and integers
$s,p\ge0$ such that
\[
 \norm{a_n b}
 \le M(1+|n|)^s\norm{(1+L_0)^p b}
\tag{3.1}\label{eq:modebound}
\]
for every finite-energy vector $b$.  Under this hypothesis, for every
$f\in C^\infty(S^1)$ the series
$Y(a,f)=\sum_n\widehat f_n a_n$ converges on finite-energy vectors,
defines a closable operator, preserves
$\Hh^\infty=\bigcap_{q\ge0}\operatorname{Dom}(L_0^q)$, and obeys
\[
 \norm{Y(a,f)b}
 \le M\,\norm{f}_{H^{s+1}}\,
       \norm{(1+L_0)^p b}.
\tag{3.2}\label{eq:smearedbound}
\]
If a unitary VOA is generated by homogeneous energy-bounded fields, all
of its fields are energy bounded.  A simple unitary VOA generated by
weight-one affine currents and conformal vectors is strongly local; its
smeared fields generate an irreducible conformal net.
\end{theorem}

This combines CKLW Proposition 6.1, Theorems 6.8, 8.1 and 8.3
\cite{CKLW}, in exactly the direction used below.  The newer extension
theorem also states that every positive-definite even lattice VOA is
energy bounded \cite[Proposition 4.2]{CarpiTomassini}.  These are
\emph{continuum} statements; they do not assert a bound uniform in a
lattice cutoff.

\begin{theorem}[Kawahigashi--Longo--M\"uger]
\label{thm:KLM}
Let $\mathcal N$ be a local irreducible conformal net satisfying the split
property and strong additivity, and suppose its two-interval inclusion
\[
 \mathcal N(E)\subset\widehat{\mathcal N}(E)
 :=\mathcal N(E')'
\]
has finite index.  Then $\mathcal N$ is completely rational, its DHR
category is a nondegenerate finite braided tensor category, and
\[
 \mu_{\mathcal N}
 =[\widehat{\mathcal N}(E):\mathcal N(E)]
 =\sum_{\rho\in\operatorname{Irr}(\mathcal N)}d(\rho)^2.
\tag{3.3}\label{eq:KLMsum}
\]
If $\mathcal N\subset\mathcal M$ is an irreducible finite-index local
extension of completely rational nets, then
\[
 \mu_{\mathcal M}
 =\frac{\mu_{\mathcal N}}{[\mathcal M:\mathcal N]^2}.
\tag{3.4}\label{eq:KLMextension}
\]
In particular $\mu_{\mathcal M}=1$ exactly when the vacuum is the only
irreducible DHR sector.
\end{theorem}

This is the multi-interval/global-index and extension theorem of
Kawahigashi--Longo--M\"uger \cite{KLM}.

\begin{theorem}[Q-systems and local extensions]
\label{thm:Qsystems}
Let $N\subset M$ be a finite-index inclusion of factors.  A conjugate of
the inclusion determines a standard Q-system $(\theta,w,x)$ in
$\operatorname{End}(N)$, and every standard Q-system reconstructs the
inclusion up to isomorphism fixing $N$.  If the endomorphisms are DHR
endomorphisms of a conformal net, the reconstructed extension is local
if and only if the Q-system is commutative with respect to the DHR
braiding.  For a finite outer group action $G\curvearrowright M$,
$N=M^G$ has index $|G|$, group-average expectation, and the regular
Q-system.
\end{theorem}

This is the Longo--Rehren reconstruction in the local-extension form
\cite{LongoRehren}.

\begin{theorem}[Trace class implies split]\label{thm:split}
Let $\mathcal N$ be a conformal net with conformal Hamiltonian $L_0$.
If
\[
 \Tr(e^{-\beta L_0})<\infty\qquad\text{for every }\beta>0,
\tag{3.5}\label{eq:traceclass}
\]
then $\mathcal N$ has the split property.  More quantitatively, the
trace-class condition implies the conformal Buchholz--Wichmann
nuclearity maps are nuclear; the Buchholz--Wichmann condition implies
statistical independence for strictly separated intervals.
\end{theorem}

The implication is due to Buchholz--Wichmann and its conformal
refinements by Buchholz--D'Antoni--Longo \cite{BuchholzWichmann,BDL}.

\begin{theorem}[Central charge eight and the $E_8$ theory]
\label{thm:DM}
Let $V$ be a simple, rational, $C_2$-cofinite, holomorphic unitary VOA of
CFT type and central charge $8$.  Then its character is
\[
 \operatorname{ch}_V(\tau)
 =\frac{E_4(\tau)}{\eta(\tau)^8}
 =q^{-1/3}(1+248q+4124q^2+34752q^3+\cdots),
\tag{3.6}\label{eq:E8char}
\]
$V_1$ is the simple Lie algebra $E_8$ at level one, and
$V\cong V_{E_8}$.  The strongly local net associated by the CKLW
construction is the level-one $E_8$ current net.
\end{theorem}

The VOA classification statement is the central-charge-eight theorem of
Dong--Mason \cite{DongMason}; the last sentence is CKLW.  We stress that
KLM alone proves holomorphy, not the name ``$E_8$''.  The VOA
reconstruction/weight-one argument is the additional identification step.

\section{L1: uniform energy bounds}
\label{sec:L1}

\subsection{The continuum estimate is complete}

\begin{lemma}[Linear energy bound for the limiting norm-two field]
\label{lem:continuumEB}
In the level-one lattice VOA of the glued rank-eight lattice, the
norm-two vector $e^\lambda$ is a weight-one affine current.  For every
$f\in C^\infty(S^1)$ there is a constant $c_f$ such that
\[
 \norm{Y(e^\lambda,f)\xi}
 \le c_f\norm{(1+L_0)\xi}.
\tag{4.1}\label{eq:continuumEB}
\]
\end{lemma}

\begin{proof}
The norm formula gives $h_\lambda=\norm{\lambda}^2/2=1$.
The weight-one subspace of the lattice VOA is the direct sum of the
Cartan currents and the root vectors $e^\alpha$ for
$\norm{\alpha}^2=2$.  Their modes satisfy the level-one affine
commutation relations.  Positivity of the Sugawara expression for $L_0$
gives, for a fixed normalized current $J$,
\[
 \norm{J_n\xi}
 \le C(1+|n|)\norm{(1+L_0)^{1/2}\xi}.
\tag{4.2}\label{eq:currentmode}
\]
Indeed, on each $L_0$ eigenspace, the positive quadratic Sugawara sum
contains $J_{-n}J_n$; commuting a negative mode to the right contributes
the central term $|n|$.  Summing
$\sum_n|\widehat f_n|\norm{J_n\xi}$ and applying Cauchy--Schwarz with
$\sum_n(1+|n|)^4|\widehat f_n|^2<\infty$ gives
\[
 \norm{J(f)\xi}
 \le C\norm{f}_{H^2}\norm{(1+L_0)\xi}.
\]
This is also the weight-one case of
Theorem~\ref{thm:CKLWenergy}.
\end{proof}

\subsection{The edge-projected lattice Green form}

We now define and estimate the concrete $G_N$ used by the collar
bosonization.  It is important that this is a definition from the frozen
QWZ edge band, not a continuum ansatz.

Use antiperiodic (NS) momenta
\[
 p_n=\frac{2\pi}{N}\left(n+\frac12\right),
 \qquad -\frac N2\le n<\frac N2,
\tag{4.3}\label{eq:NSmomenta}
\]
and Hardy occupation $\nu_n=\id_{\{n<0\}}$.  The same argument works
after a fixed deck shift because only the number of indices crossing the
Hardy cut is used.  Let $u(p)$ be the normalized top-edge eigenvector of
the $2N_y$-dimensional strip matrix through $p=0$.  Away from the
crossing choose any normalized continuation; the global bound below
uses no regularity there.

For $1\le m\le N/2$, define the finite edge-current structure factor
\[
 S_N(m)
 :=\sum_{n=-m}^{-1}
 \left|\ip{u(p_{n+m})}{u(p_n)}\right|^2,
\tag{4.4}\label{eq:structurefactor}
\]
and extend it evenly to negative $m$.  With the lattice derivative
\eqref{eq:latticesymbol}, the Green multiplier is
\[
 g_N(m):=\frac{S_N(m)}{2\omega_N(m)^2}.
\tag{4.5}\label{eq:greenmultiplier}
\]
Let
$\mathbb M_N=\{-N/2+1,\ldots,-1,1,\ldots,N/2\}$ be the fundamental
nonzero Fourier representatives.  For an eight-component zero-mean
charge $q$ in the band-limited trigonometric space $\V_N$, put
\[
 \mathcal Q_N(q):=\ip{q}{G_Nq}
 =\sum_{a=1}^8\sum_{m\in\mathbb M_N}
 g_N(m)|\widehat q_a(m)|^2.
\tag{4.6}\label{eq:greenform}
\]
The continuum chiral form is
\[
 \mathcal Q_{\rm ch}(q)
 =\frac12\sum_{a=1}^8\sum_{m\ne0}
 \frac{|\widehat q_a(m)|^2}{|m|}.
\tag{4.7}\label{eq:chiralgreen}
\]
The diagonal logarithm is omitted on both sides by the same Wick
normal ordering.

\begin{lemma}[Exact transverse edge doublet]\label{lem:edgedoublet}
At $p=0$, the $M=1$, $N_y=8$ strip Hamiltonian has two zero modes, one
on each boundary.  The remaining spectrum is $\{+1,-1\}$, each with
multiplicity seven.  The two eigenbranches through zero are analytic
near $p=0$, have slopes $+1$ and $-1$, and are separated there from the
bulk complement by $\delta_{\rm bulk}=1$.
\end{lemma}

\begin{proof}
At $p=0$ the onsite term in \eqref{eq:qwz} vanishes.  The transverse
hopping is
\[
 T_y=\frac{\sigma_y}{2i}-\frac{\sigma_z}{2}.
\]
It has rank one and singular value one.  In the eigenbasis of
$\sigma_x$, the open strip is the orthogonal sum of seven two-site
blocks
$\left(\begin{smallmatrix}0&1\\1&0\end{smallmatrix}\right)$ and two
unpaired boundary vectors.  Hence the stated spectrum and
$\delta_{\rm bulk}=1$.  The derivative of the onsite matrix at zero is
$\sigma_x$; its compression to the two boundary zero modes is diagonal
with eigenvalues $+1$ and $-1$.  Analytic perturbation theory for this
isolated two-dimensional spectral subspace, followed by diagonalization
of the first derivative inside it, gives two analytic branches with
those slopes.
\end{proof}

\begin{lemma}[Sharp lattice-symbol comparison]\label{lem:sine}
For every $N$ and $0<|m|\le N/2$,
\[
 \frac{2}{\pi}|m|\le\omega_N(m)\le |m|.
\tag{4.8}\label{eq:sinecomparison}
\]
Both constants are sharp.
\end{lemma}

\begin{proof}
Set $x=\pi|m|/N\in(0,\pi/2]$.  Concavity of $\sin x$ on this interval
places its graph above the chord joining $(0,0)$ and $(\pi/2,1)$, so
$\sin x\ge2x/\pi$.  The inequality $\sin x\le x$ gives the upper bound.
The lower constant is attained at $|m|=N/2$, and the upper constant is
approached as $|m|/N\to0$.
\end{proof}

\begin{lemma}[Particle--hole Schur bound]\label{lem:Schur}
For every $N$ and $0<|m|\le N/2$,
\[
 0\le S_N(m)\le |m|.
\tag{4.9}\label{eq:Schur}
\]
\end{lemma}

\begin{proof}
For $m>0$, exactly $m$ occupied Hardy indices
$n=-m,\ldots,-1$ are carried across the Fermi cut.  Every summand in
\eqref{eq:structurefactor} lies in $[0,1]$ by Cauchy--Schwarz and the
normalization of $u(p)$.  This proves the claim; evenness handles
$m<0$.  Equivalently, this is the Schur bound for the compression of the
momentum-shift partial isometry between occupied and unoccupied edge
subspaces.
\end{proof}

\begin{lemma}[Bulk resolvent comparison]\label{lem:resolvent}
Assume the bulk-edge decomposition \eqref{eq:bulkgap}.  For
$z$ on a contour at distance $\delta_{\rm bulk}/2$ from the bulk
spectrum and for two bulk Hamiltonians $H_{\rm b},\widetilde H_{\rm b}$,
\[
 \norm{(z-H_{\rm b})^{-1}-(z-\widetilde H_{\rm b})^{-1}}
 \le\frac{4}{\delta_{\rm bulk}^2}
       \norm{H_{\rm b}-\widetilde H_{\rm b}}.
\tag{4.10}\label{eq:resolventidentity}
\]
Consequently the corresponding bulk spectral projections differ by at
most a contour-length constant times the right-hand side.
\end{lemma}

\begin{proof}
The resolvent identity gives
\[
 (z-H_{\rm b})^{-1}-(z-\widetilde H_{\rm b})^{-1}
 =(z-H_{\rm b})^{-1}(H_{\rm b}-\widetilde H_{\rm b})
   (z-\widetilde H_{\rm b})^{-1}.
\]
Both resolvent norms are at most $2/\delta_{\rm bulk}$.  Integrating the
identity in the Riesz formula proves the projection statement.
\end{proof}

\begin{remark}
Lemma~\ref{lem:resolvent} controls the massive transverse complement.
It supplies the isolated analytic edge subspace used in
Lemma~\ref{lem:edgedoublet}, but it cannot by itself furnish the
$|m|^{-1}$ bound in \eqref{eq:chiralgreen}.  That bound comes from the
exact Hardy particle--hole count and Lemma~\ref{lem:sine}.
\end{remark}

\begin{theorem}[Uniform Green-form domination and Sobolev convergence]
\label{thm:UGF}
For the edge-projected frozen $M=1$, $N_y=8$ QWZ collar, every
zero-mean $q\in\V_N\otimes\C^8$ satisfies
\[
 0\le\mathcal Q_N(q)
 \le\frac{\pi^2}{4}\mathcal Q_{\rm ch}(q).
\tag{4.11}\label{eq:UGFhere}
\]
Thus one may take the all-volume constant
\[
 K_G=\frac{\pi^2}{4}=2.467401100272\ldots .
\tag{4.12}\label{eq:KG}
\]
If $P_N$ is Fourier projection onto $|m|\le N/2$, then
$G_NP_N\to G_{\rm ch}$ in local Sobolev quadratic-form norm:
\[
 \left|\mathcal Q_N(P_Nq)-\mathcal Q_{\rm ch}(q)\right|
 \le\varepsilon_N\norm{q}_{H^2(S^1;\C^8)}^2,
 \qquad\varepsilon_N\longrightarrow0.
\tag{4.13}\label{eq:SobolevUGF}
\]
\end{theorem}

\begin{proof}
Lemmas~\ref{lem:sine} and \ref{lem:Schur} give, mode by mode,
\[
 0\le g_N(m)
 =\frac{S_N(m)}{2\omega_N(m)^2}
 \le\frac{|m|}{2(2|m|/\pi)^2}
 =\frac{\pi^2}{8|m|}
 =\frac{\pi^2}{4}\frac1{2|m|}.
\tag{4.14}\label{eq:modeUGF}
\]
Multiplication by $|\widehat q_a(m)|^2$ and summation over the modes and
eight independent components proves \eqref{eq:UGFhere}.  Notice that
the number of components does not multiply $K_G$.

It remains to prove convergence.  Fix $m$.  The $|m|$ terms in
\eqref{eq:structurefactor} involve momenta satisfying
$|p_n|+|p_{n+m}|=O_m(N^{-1})$.  By
Lemma~\ref{lem:edgedoublet}, $u(p)$ is norm-continuous there, uniformly
over this finite set.  Consequently
\[
 S_N(m)\longrightarrow |m|,
\qquad
 \frac{\omega_N(m)}{|m|}\longrightarrow1,
\qquad
 g_N(m)\longrightarrow\frac1{2|m|}.
\tag{4.15}\label{eq:pointwisegreen}
\]
Define
\[
 \varepsilon_N:=
 \sup_{0<|m|\le N/2}
 \frac{|g_N(m)-(2|m|)^{-1}|}{(1+m^2)^2}
 +\sup_{|m|>N/2}\frac{1}{2|m|(1+m^2)^2}.
\tag{4.16}\label{eq:epsilonN}
\]
For any fixed $M$, the first supremum over $|m|\le M$ tends to zero by
\eqref{eq:pointwisegreen}.  For $|m|>M$,
\eqref{eq:modeUGF} bounds the numerator by
$(K_G+1)/(2|m|)$, uniformly in $N$; hence the remaining supremum tends
to zero as $M\to\infty$.  Thus $\varepsilon_N\to0$.  Parseval's identity
and the definition of the $H^2$ norm now give
\eqref{eq:SobolevUGF}.
\end{proof}

\begin{theorem}[L1 after the proved Green estimate]
\label{thm:L1reduction}
For the edge-projected bosonization, assume only that the normalized
cocycles and zero-mode charge shifts have norm one in all four grades.
Then for every smooth $f$ the bound \eqref{eq:L1} holds with $p=1$ and
\[
 C_f\le C(\pi^2/4,\lambda)\norm{f}_{H^2},
\tag{4.17}\label{eq:Cf}
\]
independent of $N$.
\end{theorem}

\begin{proof}
Bosonic Wick ordering factors a mode of \eqref{eq:chargedfield} into the
norm-one zero-mode shift and exponentials of creation and annihilation
currents.  The normal-ordered lattice-VOA estimate is the finite-mode
version of Theorem~\ref{thm:CKLWenergy}; its constant depends on the
Gaussian contraction only through the Green quadratic form.  Theorem
\ref{thm:UGF} bounds that contraction uniformly by the continuum form
with $K_G=\pi^2/4$.  The norm-two vector has weight one, so
Lemma~\ref{lem:continuumEB} and \eqref{eq:smearedbound} give degree
$p=1$ and \eqref{eq:Cf}.  Wick scalar factors are bounded by the same
quadratic-form domination.
\end{proof}

\begin{remark}[Scope of the proved UGF]
Theorem~\ref{thm:UGF} is an all-$N$ theorem for the edge-projected
bosonization used in \eqref{eq:chargedfield}; it is not a fit to the six
tested sizes.  It does not prove that every microscopic full-CAR density
operator belongs to this bosonic edge algebra.  Excluding additional
bulk/ultraviolet local scaling operators is exactly the separate
\textup{(ALG-EXH)} condition.  A completely gapless transverse covariance
would contribute $O(N)$ particle--hole channels already at $m=1$ and
would violate \eqref{eq:UGFhere}; the isolated one-branch hypothesis is
therefore load-bearing.
\end{remark}

\paragraph{Finite audit (numerical support, not proof).}
The companion probe
\path{experiments/tfpt-discovery/mmst_ugf_bound_probe.py}
checks every displayed inequality at
$N=16,24,32,48,64,96$.  It finds
$\max_{N,m}2|m|g_N(m)=1.029399027646<K_G$, verifies
$\delta_{\rm bulk}=1$ and the unit edge slope, and rejects both a
gapless transverse mutant (whose ratio grows from $64.829$ to $384.137$)
and the false symbol constant $\omega_N(m)\ge|m|$.

\section{L2: strong convergence and local algebras}
\label{sec:L2}

\subsection{Explicit dyadic mode and Fock isometries}

Let
\[
 \mathcal K_N
 =\ell^2\!\left(\{m\in\Z:0<|m|\le N/2\}\right)\otimes\C^8
\tag{5.1}\label{eq:oneparticleboson}
\]
with the single Nyquist representative convention of
$\mathbb M_N$.  Define
\[
 j_{N,2N}(e_m\otimes e_a)=e_m\otimes e_a,
\qquad m\in\mathbb M_N,\quad 1\le a\le8.
\tag{5.2}\label{eq:modeisometry}
\]
Thus $j_{N,2N}$ is the zero-extension/Shannon embedding of the old
Fourier modes into the refined lattice.  It is an isometry and
\[
 j_{2N,4N}j_{N,2N}=j_{N,4N}.
\tag{5.3}\label{eq:modecompat}
\]
On the charge-lattice sum of symmetric bosonic Fock spaces set
\[
 V_{N,2N}
 :=\id_{\ell^2(L_{\rm charge})}\otimes\Gamma_s(j_{N,2N}).
\tag{5.4}\label{eq:Fockisometry}
\]
Second quantization preserves isometries and composition, so
$V_{N,2N}^*V_{N,2N}=1$ and
$V_{2N,4N}V_{N,2N}=V_{N,4N}$.  The charge-lattice identity makes the
same statement valid in every one of the four $\lambda$ grades.

\subsection{The multiplier rate extracted from UGF}

For $m>0$ define the vertex multiplier
\[
 r_N(m)
 :=\frac{\sqrt{mS_N(m)}}{\omega_N(m)}
 =\sqrt{2m\,g_N(m)}
\quad (m\le N/2),
\qquad
 r_N(m):=0\quad(m>N/2).
\tag{5.5}\label{eq:vertexmultiplier}
\]
The creation part of the finite vertex exponential contains
$r_N(m)a_{-m}/m$; the annihilation part is its adjoint.  Theorem
\ref{thm:UGF} gives $0\le r_N(m)\le\pi/2$.

\begin{lemma}[Uniform quadratic multiplier error]\label{lem:multiplier}
For all even $N\ge16$ and $1\le m\le N/2$,
\[
 |r_N(m)-1|
 \le A_0\frac{m^2}{N^2},
\qquad
 A_0:=4096\left(1+\frac\pi2\right).
\tag{5.6}\label{eq:rNrate}
\]
Consequently, for $1\le m\le N$,
\[
 |r_{2N}(m)-r_N(m)|
 \le A_{\rm TEL}\frac{m^2}{N^2},
\qquad
 A_{\rm TEL}:=\frac54A_0
 =5120\left(1+\frac\pi2\right).
\tag{5.7}\label{eq:dyadicmultiplier}
\]
\end{lemma}

\begin{proof}
Put $t=m/N$.  For $t\le1/64$, every momentum entering
\eqref{eq:structurefactor} lies in $[-\pi/32,\pi/32]$.  Differentiating
the explicit $M=1$ transfer recurrence gives
\[
 \sup_{|p|\le\pi/32}\norm{u'(p)}\le4.
\tag{5.8}\label{eq:edgederivative}
\]
For completeness, the decaying transfer root is
$\rho(p)=1-\cos p$; on this interval
$|\rho|\le1-\cos(\pi/32)$ and $|\rho'|\le\sin(\pi/32)$.
Differentiating the normalized geometric vector gives a norm below
$1$, while the opposite-boundary correction is bounded by
$2|\rho|^{N_y}/(1-|\rho|)$ and its derivative by the corresponding
geometric derivative; their sum is below $4$ for $N_y=8$.
Therefore
\[
 0\le1-\frac{S_N(m)}m
 \le\norm{u(p_{n+m})-u(p_n)}^2
 \le64\pi^2t^2.
\tag{5.9}\label{eq:Srate}
\]
Also, elementary Taylor remainder and
Lemma~\ref{lem:sine} give
\[
 0\le\frac{\pi t}{\sin(\pi t)}-1\le\pi^2t^2
 \qquad(0\le t\le1/2).
\tag{5.10}\label{eq:symbolrate}
\]
Since
$r_N=(\pi t/\sin\pi t)\sqrt{S_N/m}$, these estimates prove
\eqref{eq:rNrate} for $t\le1/64$ with a constant much smaller than
$A_0$.  For $t\ge1/64$, the UGF bound gives
$|r_N-1|\le1+\pi/2\le A_0t^2$, proving the global statement.

For $m\le N/2$, apply \eqref{eq:rNrate} at $N$ and $2N$ and use the
triangle inequality, giving the factor $1+1/4=5/4$.  For
$N/2<m\le N$, $r_N(m)=0$ and
$r_{2N}(m)\le\pi/2\le2\pi m^2/N^2<A_{\rm TEL}m^2/N^2$.
\end{proof}

\begin{lemma}[Lipschitz Wick estimate]\label{lem:WickLipschitz}
Let $r,s$ be two cutoff multiplier sequences bounded by $\pi/2$ and
suppose $|r(m)-s(m)|\le a m^2$.  For the norm-two charge
$\lambda$ and every $f\in C^\infty(S^1)$,
\[
 \norm{\bigl(Y_r(e^\lambda,f)-Y_s(e^\lambda,f)\bigr)
        (1+L_0)^{-3}}
 \le C_\lambda\,a\,\norm{f}_{H^4},
\tag{5.11}\label{eq:WickLipschitz}
\]
where $C_\lambda$ is independent of both cutoffs.  The same estimate
holds for the adjoints.
\end{lemma}

\begin{proof}
Interpolate $r_\tau=(1-\tau)r+\tau s$ and differentiate the finite
normal-ordered exponential.  Every differentiated monomial has exactly
one insertion
\[
 J_{s-r}(z)
 =\sum_{m>0}\frac{s(m)-r(m)}m
   \bigl(a_{-m}z^m-a_mz^{-m}\bigr)
\]
and otherwise the vertex exponential with multiplier $r_\tau$.
The oscillator inequalities
$\norm{a_{\pm m}(1+L_0)^{-1/2}}\le\sqrt{m+1}$, followed by
Cauchy--Schwarz in $m$ and four integrations by parts in the smearing
variable, give
\[
 \norm{:\!J_{s-r}Y_{r_\tau}(e^\lambda)\!:(f)
       (1+L_0)^{-3}}
 \le C_\lambda a\norm{f}_{H^4}.
\]
The constant is uniform because $0\le r_\tau(m)\le\pi/2$, so all
undifferentiated Wick contractions are dominated by
Theorem~\ref{thm:UGF}.  Integrating over $\tau\in[0,1]$ proves
\eqref{eq:WickLipschitz}.  The finite-cutoff calculation is an
absolutely convergent sum after the displayed graph weight; monotone
cutoff removal gives the asserted bound.  Replacing $\lambda$ by
$-\lambda$ and conjugating the coefficients proves the adjoint
statement.
\end{proof}

\begin{theorem}[TEL-a: dyadic charged-field graph defect]
\label{thm:TELa}
The explicit isometries \eqref{eq:Fockisometry} satisfy, for every
$f\in C^\infty(S^1)$,
\[
\begin{aligned}
 &\norm{\bigl(\Psi_{\lambda,2N}(f)V_{N,2N}
       -V_{N,2N}\Psi_{\lambda,N}(f)\bigr)(1+L_{0,N})^{-3}}\\
 &\quad+
 \norm{\bigl(\Psi_{\lambda,2N}(\overline f)^*V_{N,2N}
       -V_{N,2N}\Psi_{\lambda,N}(\overline f)^*\bigr)
       (1+L_{0,N})^{-3}}\\
 &\hspace{35mm}\le
 2C_\lambda A_{\rm TEL}\norm{f}_{H^4}N^{-2}.
\end{aligned}
\tag{5.12}\label{eq:TELa}
\]
Thus TEL-a holds with $p=3$ and $\epsilon=1$.
\end{theorem}

\begin{proof}
After applying $V_{N,2N}$, the coarse field is the fine Fock-space
vertex field with multiplier $r_N$ extended by zero on the new modes.
Apply Lemma~\ref{lem:WickLipschitz} with
$a=A_{\rm TEL}N^{-2}$ from Lemma~\ref{lem:multiplier}; apply it again
to the adjoint.  This proves \eqref{eq:TELa}.
\end{proof}

\begin{lemma}[Summable graph defects imply strong convergence]
\label{lem:telescoping}
The dyadic estimate \eqref{eq:TELa} implies that the transported fields
and adjoints converge on $\Dfin$ along every dyadic cofinal sequence.
\end{lemma}

\begin{proof}
Let $\xi$ have energy at most $E$ and represent it at scale $N_0$.
For $K>L$, insert the compatible isometries and telescope:
\[
\norm{\Psi_{\lambda,2^KN_0}^{\rm tr}(f)\xi
      -\Psi_{\lambda,2^LN_0}^{\rm tr}(f)\xi}
\le
2C_\lambda A_{\rm TEL}\norm{f}_{H^4}(1+E)^3
\sum_{j=L}^{K-1}(2^jN_0)^{-2}.
\tag{5.13}\label{eq:telestimate}
\]
The geometric tail tends to zero.  The adjoint half of
\eqref{eq:TELa} gives graph convergence, and L1 makes the limit
closable.  A dyadic subsystem is cofinal in the scaling parameter, so
this is sufficient for the edge-mode scaling limit.
\end{proof}

\paragraph{Finite audit (numerical support, not proof).}
The companion probe
\path{experiments/tfpt-discovery/mmst_tel_probe.py}
checks \eqref{eq:modeisometry}--\eqref{eq:dyadicmultiplier} at
$N=16,24,32,48,64,96$.  The Gram and compatibility errors are exactly
zero, the measured coefficient envelope is
$\max N^2|r_{2N}(m)-r_N(m)|/m^2=3.420384012310$, and a smooth-test Wick
seminorm fits $N^{-2.022861}$.  A wrong-rate $1/N$ mutant and a
non-isometric embedding mutant are both rejected.

\begin{remark}[Local support]
The Shannon embedding \eqref{eq:modeisometry} is exact in momentum and
is not sharply support-preserving in position.  Theorem~\ref{thm:TELa}
therefore proves graph convergence of the edge-projected fields, not by
itself convergence of sharply localized microscopic algebras.  The
support/locality statement remains explicitly in \textup{(ALG-EXH)}.
\end{remark}

\subsection{TEL-b: the two-edge full-CAR remainder}

The exact Hardy model gives the quarter-twist coefficient
\[
 \norm{P_0-P_{1/4}}_{\HS}^2
 =\frac{2}{\pi^2}\log N+O(1),
\qquad \frac2{\pi^2}=0.202642367\ldots .
\tag{5.14}\label{eq:hardylog}
\]
For the frozen QWZ exhaustion, the computed decomposition is
\[
 C_N^{(0)}-C_N^{(1)}
 =D_{N,\mathrm{top}}^{(1/4)}
  \oplus D_{N,\mathrm{bottom}}^{(1/4)}+R_N,
\tag{5.15}\label{eq:tworemainder}
\]
with
\[
 \max_{N\in\{16,24,32,48,64,96\}}\norm{R_N}_{\HS}
 =2.105079358<2.11.
\tag{5.16}\label{eq:measured211}
\]
A $1/N$ fit predicts $2.107565497$, with maximum fit residual
$2.349\times10^{-5}$.  These values support boundedness of $R_N$ but do
not prove either
\[
 \sup_N\norm{R_N}_{\HS}\le2.11
\tag{5.17}\label{eq:allN211}
\]
for the full-CAR remainder.

\begin{proposition}[Exact Hankel reduction of TEL-b]
\label{prop:HankelReduction}
Let
\[
 C(p)=\id_{(-\infty,0)}(h_{\rm strip}(p))
       +\frac12\id_{\{0\}}(h_{\rm strip}(p))
\]
be the explicit $16\times16$ frozen QWZ Fermi symbol and let
$P_{\rm t},P_{\rm b}$ be the rank-one boundary projections at $p=0$.
The symbol left after the constant-profile Hardy subtraction is
\[
 B(p)=C(p)
 -\id_{(0,\pi)}(p)P_{\rm t}
 -\id_{(-\pi,0)}(p)P_{\rm b}.
\tag{5.18}\label{eq:regularizedsymbol}
\]
Writing $B(p)=\sum_{k\in\Z}B_ke^{ikp}$, every finite Hankel corner
obeys
\[
 \norm{H_N(B)}_{\HS}^2
 \le\sum_{k\ne0}|k|\norm{B_k}_{\mathrm F}^2.
\tag{5.19}\label{eq:PellerHS}
\]
Hence an all-$N$ TEL-b proof by this subtraction would follow from
\[
 16\sum_{k\ne0}|k|\norm{B_k}_{\mathrm F}^2\le C_R^2.
\tag{5.20}\label{eq:Hankeltarget}
\]
\end{proposition}

\begin{proof}
In the untwisted Hardy basis, multiplication by $B$ has matrix entries
$B_{j-k}$.  Its off-diagonal Hardy corner has entries $B_{j+k+1}$.
Counting the pairs $(j,k)\in\mathbb N_0^2$ with $j+k+1=\ell$ gives the exact
Hilbert--Schmidt identity
\[
 \norm{H(B)}_{\HS}^2
 =\sum_{\ell\ge1}\ell\norm{B_\ell}_{\mathrm F}^2.
\]
Finite corners are compressions and cannot increase this norm.  The
opposite edge uses the negative coefficients.  There are at most four
corner/orientation contributions; the triangle inequality gives the
safe factor sixteen after squaring in \eqref{eq:Hankeltarget}.  This is
the $p=2$ case of the Peller $B^{1/2}_{2,2}$ criterion \cite{Peller}.
\end{proof}

\begin{remark}[Second-angle verdict: \texttt{TEL-B-EXTERNAL}]
The new Hankel route does not certify the requested constant.  On a
$4096$-point symbol grid the companion probe finds
\[
 \sum_{k\ne0}|k|\norm{B_k}_{\mathrm F}^2
 =4.482956087,\qquad
 4\norm{B}_{\dot H^{1/2}}=8.469196974>3.
\]
The constant edge-profile subtraction cancels the infrared jump at
$p=0$ but exposes a Brillouin-seam ultraviolet contribution; the raw
symbol seminorm on the same grid is $2.570232705$.  Thus neither a
$C_R=2.5$ nor a $C_R=3$ proof follows, and the value is numerical rather
than an all-mode enclosure.  UGF cannot repair this because it controls
only the scalar edge branch, whereas \eqref{eq:regularizedsymbol} is a
matrix-valued transverse cancellation.  This was the second distinct
in-house attempt; by the stop rule TEL-b is now typed
\texttt{TEL-B-EXTERNAL}.  No all-$N$ constant is claimed.
\end{remark}

\begin{remark}[Round-3 supersession, 2026-09-05]
The preceding verdict records the historical raw-Hankel attempt.  It is not the
current TEL-b residue.  Native $256$-bit Arb/Acb arithmetic on $5812$ rational
cells now proves
\[
 \mathrm{TV}_p(F'')\le15.8176959629889243576\ldots<15.928674,
 \qquad K_3^{\mathrm{TV}}<2.517464500835,
\]
including the exact Hilbert--Schmidt cut jump $2$
(\texttt{v1022\_telb\_round3\_certificates.py}).  Separately, a $100$-bit
certificate over $3888$ analytic rectangles proves
$C_\star<0.1525$, and the Peano/Euler argument proves
$\sup_k|\varphi_N(k)|<0.509/N$ for every even $N\ge16$
(\texttt{v1025\_telb\_c1\_c2a\_certificates.py}).  These close the TV, C1 and
C2a inputs. Round 4 additionally proves CF $<0.065697$, DG $<0.042959$,
and the complete relaxed norm $\norm{R_N}_{\HS}<2.995906<3$ for every even
$N\ge16$, including the finite-grid aliases
(\texttt{v1026\_telb\_round4\_closure.py}). The remaining microscopic
algebra identification is not a consequence of this norm theorem.
\end{remark}

\subsection{Third angle: certified finite part plus a dyadic tail}

There are two natural full-CAR one-particle embeddings from circumference
$N$ to $2N$.  The low-energy embedding $J_N^{\rm low}$ preserves the
integer Fourier label $n$, while the physical-momentum dilation
$J_N^{\rm dil}$ sends $n\mapsto2n$.  The proposed tail mechanism would
require, for one of them,
\[
 \norm{R_{2N}-J_NR_NJ_N^*}_{\HS}\le DN^{-2}
\tag{5.21}\label{eq:Rtailproposed}
\]
with one explicit $D$.  The reverse triangle inequality would then give
\[
 \left|\norm{R_{2N}}_{\HS}-\norm{R_N}_{\HS}\right|
 \le DN^{-2},
\tag{5.22}\label{eq:Rnormtail}
\]
and the dyadic tail from $N_0$ would be at most
$DN_0^{-2}/(1-1/4)$.

\begin{proposition}[The TEL-a embedding does not supply the proposed
remainder tail]\label{prop:tailfailure}
On the frozen finite matrices, neither natural embedding exhibits the
defect required by \eqref{eq:Rtailproposed}.  The measured values are
\[
\begin{array}{c|c|c}
N\to2N&
\norm{R_{2N}-J_N^{\rm low}R_N(J_N^{\rm low})^*}_{\HS}&
\norm{R_{2N}-J_N^{\rm dil}R_N(J_N^{\rm dil})^*}_{\HS}\\ \hline
16\to32&2.071346866&2.066044074\\
32\to64&2.067361892&2.072036965\\
64\to128&2.065611659&2.075227063\\
96\to192&2.065062404&2.076335188
\end{array}
\tag{5.23}\label{eq:Rdefecttable}
\]
Thus the certified-tail argument cannot inherit the proved TEL-a
$N^{-2}$ comparison.
\end{proposition}

\begin{proof}
In the untwisted Fourier basis, construct $R_N$ exactly as in
\eqref{eq:tworemainder}; both embeddings are coordinate isometries.
For either $J$,
\[
 \norm{R_{2N}-JR_NJ^*}_{\HS}^2
 =\norm{R_{2N}}_{\HS}^2+\norm{R_N}_{\HS}^2
 -2\Re\ip{J^*R_{2N}J}{R_N}_{\HS}.
\]
The companion probe evaluates this identity without forming $J$ and
reproduces the established finite norms.  The table is numerical
evidence, not an interval theorem; it shows that the proposed small
TEL-a-type defect is absent on the tested scales.  Therefore
high-precision recalculation of the same finite values cannot certify
the missing infinite tail.
\end{proof}

\begin{remark}[The strictly smaller dyadic residue]
The scalar norm increments are instead
\[
 0.007486497094,\quad0.003790483818,\quad
 0.001913049365,\quad0.001279839965
\]
for the four rows of \eqref{eq:Rdefecttable}; the products $N$ times
these increments are
\[
 0.119783953,\quad0.121295482,\quad
 0.122435159,\quad0.122864637.
\]
This suggests, but does not prove,
\[
 \left|\norm{R_{2N}}_{\HS}-\norm{R_N}_{\HS}\right|
 \le\frac{A_R}{N}\qquad(N\ge96).
\tag{5.24}\label{eq:Rscalarresidue}
\]
Unlike an adjacent-size $1/N$ estimate, \eqref{eq:Rscalarresidue} is
summable on a dyadic tower.  If proved, it gives
\[
 \sup_{k\ge0}\norm{R_{96\cdot2^k}}_{\HS}
 \le\norm{R_{96}}_{\HS}+\frac{2A_R}{96}.
\tag{5.25}\label{eq:Rdyadicbound}
\]
The edge-mode scaling theorem uses a cofinal dyadic subsystem, so
\eqref{eq:Rdyadicbound}, not a supremum over every integer $N$, is the
minimal required statement.  The six reported finite values are
ordinary double-precision diagnostics.  Arbitrary-precision
\texttt{mpmath} output alone is not a rigorous interval certificate;
because the prerequisite tail theorem failed, no floating-point value
is relabelled as certified.
\end{remark}

\begin{remark}[The dyadic-increment route is retired]
Equations~\eqref{eq:Rscalarresidue}--\eqref{eq:Rdyadicbound} are retained as
an audit trail, not as the current proof route.  The cover/sawtooth reduction
and the round-3 C1/C2a/TV certificates supersede this proposed scalar
increment.  No all-$N$ conclusion is imported from the measurements above.
\end{remark}

\paragraph{Third-angle audit.}
\path{experiments/tfpt-discovery/mmst_telb_tail_probe.py} reproduces all
six established finite norms, evaluates both embeddings through
$N=192$, verifies the reverse-triangle inequality, and rejects an
$N^{-1}$ mutant as an $N^{-2}$ tail.  Its verdict is
\texttt{TAIL-TELESCOPING-FAILS}; all diagnostic and honesty guards pass.

\subsection{Relative locality of the charged and neutral fields}

\begin{theorem}[Relative locality survives the graph limit]
\label{thm:relativeLocality}
The limiting edge-projected charged field is relatively local with
respect to the neutral $D_5\oplus A_3$ current net.
\end{theorem}

\begin{proof}
For every neutral root $\alpha$, the frozen generator satisfies
\[
 \ip{\lambda}{\alpha}\in\Z.
\tag{5.26}\label{eq:integralpairing}
\]
On the displayed simple-root basis the eight pairings are
$[0,0,0,0,1,1,0,0]$; integrality for the full root lattice follows by
linearity.  Therefore the finite vertex exchange phase is
$e^{2\pi i\ip{\lambda}{\alpha}}=1$.  If $f$ and $g$ have disjoint
supports, the finite collar consequently satisfies
\[
 [\Psi_{\lambda,N}(f),W_{\alpha,N}(g)]\xi=0
\]
on the finite-energy core, where $W_{\alpha,N}(g)$ is a bounded neutral
Weyl generator.

Transport to the common Hilbert space with \eqref{eq:Fockisometry}.
Theorem~\ref{thm:TELa} gives graph convergence of the charged field and
its adjoint; Osborne--Stottmeister gives strong convergence of the
bounded neutral Weyl operators.  For a common-core vector $\xi$,
\[
\begin{aligned}
 W_\alpha(g)\Psi_\lambda(f)\xi
 &=\lim_NW_{\alpha,N}(g)\Psi_{\lambda,N}(f)\xi\\
 &=\lim_N\Psi_{\lambda,N}(f)W_{\alpha,N}(g)\xi
 =\Psi_\lambda(f)W_\alpha(g)\xi .
\end{aligned}
\]
The uniform graph bound justifies both products.  Closability extends
the equality from the core, and the Weyl generators strongly generate
the neutral local algebra.
\end{proof}

\subsection{Outerness of the limiting four-grade action}

\begin{theorem}[Outerness from the connected simple-current Q-system]
\label{thm:outerness}
On every limiting local factor generated by the neutral algebra and the
four charged grades, every nonidentity element of the $\Z_4$ gauge
action is outer.
\end{theorem}

\begin{proof}
Let $\iota:\A(I)\hookrightarrow\B^{\rm gen}(I)$ be the local inclusion.
Its canonical endomorphism has the multiplicity-free decomposition
\[
 \overline\iota\iota
 \simeq\id\oplus\rho_\lambda\oplus\rho_{2\lambda}
              \oplus\rho_{3\lambda}.
\tag{5.27}\label{eq:canonicalendo}
\]
The four sectors are inequivalent: their exact grades are distinct and
$[\lambda]$ has order four.  Only the first summand is the vacuum, so
\[
 \A(I)'\cap\B^{\rm gen}(I)
 \cong\operatorname{Hom}(\id,\overline\iota\iota)=\C.
\tag{5.28}\label{eq:irreducibleinclusion}
\]
Suppose $\gamma_g=\operatorname{Ad}u$ for a nontrivial $g\in\Z_4$ and
$u\in\B^{\rm gen}(I)$.  Since $\gamma_g$ fixes $\A(I)$ pointwise,
$u\in\A(I)'\cap\B^{\rm gen}(I)=\C$, so $\gamma_g=\id$.  This contradicts
faithfulness: on the grade-one field it acts by the nontrivial character
$\gamma_g(\Psi_\lambda)=\chi_1(g)\Psi_\lambda$.  Hence $\gamma_g$ is
outer.

There is no finite-dimensional outerness witness: every automorphism of
$M_{256}$ is inner, and
$U=\operatorname{diag}(1_{64},i1_{64},-1_{64},-i1_{64})$ is a
charge-zero implementing unitary.  Outerness is a local-factor limit
statement derived from \eqref{eq:irreducibleinclusion}, not inherited
from finite matrix algebras.
\end{proof}

\subsection{From fields to local algebras}

For an interval $I$, define the field-generated finite algebra
\[
 \B_N^{\rm gen}(I)
 =W^*\!\left(
 \A_N(I),\
 \overline{\Psi_{k\lambda,N}(f)}:
 k\in\Z_4,\ \supp f\subset I
 \right).
\tag{5.29}\label{eq:generatedfinite}
\]
Define $\B^{\rm gen}(I)$ similarly from the limiting fields.

\begin{lemma}[Graph convergence of a complete generating family]
\label{lem:algconv}
Assume L1, the graph convergence \eqref{eq:graphconv} for a
support-preserving countable dense family of test functions, strong
locality of the limiting fields, and the exhaustion condition
\textup{(ALG-EXH)} from the final box.  Then
\[
 \B_N^{\rm gen}(I)\longrightarrow\B^{\rm gen}(I)
\tag{5.30}\label{eq:algconv}
\]
in the standard scaling-algebra sense: bounded functions of the
generating fields converge strongly, the strong lower and upper limits
coincide, and no additional local scaling operator remains.
\end{lemma}

\begin{proof}
Graph convergence of a closed operator and its adjoint implies strong
resolvent convergence for its hermitian real and imaginary parts on the
common invariant core.  Bounded continuous functions of those parts
therefore converge strongly.  Their spectral projections generate the
same von Neumann algebra as the affiliated smeared fields.  Thus every
element of the limiting generated algebra belongs to the strong lower
limit.  Conversely, \textup{(ALG-EXH)} says exactly that every bounded local
scaling sequence is approximable by bounded words in this generating
family; its weak cluster points lie in $\B^{\rm gen}(I)$.  Hence the
upper limit is contained in the lower limit.
\end{proof}

\begin{proposition}[Sharp Buchholz--Verch exhaustion reduction]
\label{prop:BVexhaustion}
Let $\underline{\B}^{\rm gen}(I)$ denote the scaling subalgebra generated
by the neutral Weyl operators and the four graph-convergent charged
grades.  Equality with the full microscopic Buchholz--Verch scaling
algebra follows if the following two estimates hold:
\[
\begin{split}
 &\lim_{E\to\infty}\ \sup_{N\ge16}\
 \sup_{\substack{x\in\B_N(I)\\\norm{x}\le1}}
 \norm{(1-P^{\rm gen}_{E,N})e^{-\beta L_{0,N}}x\Omega_N}=0
 \quad(\beta>0),\\
 &\lim_{N\to\infty}\ 
 \inf_{y_N\in\B_N^{\rm gen}(I)}
 \left(
 \norm{(x_N-y_N)P_{E,N}}
 +\norm{(x_N^*-y_N^*)P_{E,N}}
 \right)=0
\end{split}
\tag{5.31}\label{eq:EXHsharp}
\]
for every finite $E$ and every norm-bounded, translation-continuous
local scaling sequence $(x_N)$.
\end{proposition}

\begin{proof}
The first line is uniform phase-space tightness: after Euclidean energy
damping, vectors created by the microscopic unit ball are uniformly
approximated by the finite-energy neutral-plus-four-grade subspace.  The
second line says that no bounded operator sequence has additional
matrix elements on any finite-energy sector, for the operator and its
adjoint.  It gives equality of the strong lower limits on the common
finite-energy core.  Phase-space tightness then excludes a weak cluster
point supported only at energies escaping to infinity, giving equality
of the upper limit as well.  This is the standard scaling-algebra
compactness mechanism of Buchholz--Verch \cite{BuchholzVerch}, written
here as the two concrete estimates the collar must satisfy.
\end{proof}

\begin{remark}[Verified and missing Buchholz--Verch hypotheses]
The frozen lattice system has isotony, exact finite-range locality,
uniform norm bounds for scaling families, discrete translation
covariance, and a translation-covariant neutral scaling limit.
Theorems~\ref{thm:UGF}, \ref{thm:TELa} and
\ref{thm:relativeLocality} add uniform charged energy control, graph
convergence and the lower inclusion
$\underline{\B}^{\rm gen}(I)\subset\underline{\B}(I)$.  What is not
proved is the uniform microscopic phase-space estimate and the second,
upper-inclusion line of \eqref{eq:EXHsharp}.  These are precisely the
``no extra ultraviolet operators'' content; no arithmetic or finite
sector census can establish them.
\end{remark}

\begin{remark}[Binding ALG2 correction, 2026-09-05]
The former claim that Fourier/Vandermonde interpolation automatically proves
the second line of \eqref{eq:EXHsharp} for
$\B_N^{\rm gen}(I)$ is withdrawn.  The exact one-mode countermodel in
\texttt{v1022\_telb\_round3\_certificates.py} has full CAR algebra $M_2$,
whereas its even/code-generated subalgebra is diagonal; the annihilator is at
distance at least $1$ from that subalgebra.  Vandermonde inversion reaches a
full-CAR compression, but does not identify the microscopic generated algebra.
Consequently the upper inclusion requires an independent FE-GEN
compression-surjectivity theorem, or a precise gauge/code-projected
microscopic-algebra argument.  The positive phase-space/ALG1 reading is also
conditional on the generated projection containing the microscopic
low-energy sectors.
\end{remark}

\begin{theorem}[Exact reduction of L2]\label{thm:L2reduction}
Assume the exhaustion estimates \eqref{eq:EXHsharp}.  Then L2 of
Theorem~\ref{thm:master} holds.
\end{theorem}

\begin{proof}
Lemma~\ref{lem:telescoping} gives the nonzero charged field after fixing
one coherent cocycle/phase normalization.
Theorem~\ref{thm:relativeLocality} and the integer spin $h_\lambda=1$
make the order-four Q-system commutative.
The neutral stress tensor converges by Theorem~\ref{thm:OS}; covariance
of the charged primary fields then supplies conformal covariance of the
generated extension.  Lemma~\ref{lem:algconv} gives local-algebra
convergence, using Proposition~\ref{prop:BVexhaustion}.
Theorem~\ref{thm:outerness} supplies the outer local action.
Nonzeroness follows from the normalized vacuum-to-charge
matrix element fixed to one at every scale and preserved by the graph
limit.
\end{proof}

\section{\texorpdfstring{L3: finite-group $C^*$ crossed products commute with the limit}{L3: finite-group crossed products commute with the limit}}
\label{sec:L3}

\subsection{\texorpdfstring{The $C^*$-inductive theorem}{The C-star inductive theorem}}

\begin{theorem}[Finite-group crossed-product/limit exchange]
\label{thm:crossedlimit}
Let $G$ be finite and let
\[
 A_1\xrightarrow{\iota_1}A_2\xrightarrow{\iota_2}\cdots
\]
be an inductive system of unital $C^*$-algebras with injective maps.
Suppose $\alpha^{(N)}:G\curvearrowright A_N$ are actions and
\[
 \iota_N\alpha_g^{(N)}
 =\alpha_g^{(N+1)}\iota_N .
\tag{6.1}\label{eq:equivariance}
\]
Let $A=\varinjlim A_N$ and $\alpha=\varinjlim\alpha^{(N)}$.  Then the
canonical map is a $*$-isomorphism
\[
 \varinjlim(A_N\rtimes_{\alpha^{(N)}}G)
 \cong A\rtimes_\alpha G.
\tag{6.2}\label{eq:crossediso}
\]
For the coefficient expectations
\[
 F_N\!\left(\sum_{g\in G}a_gu_g\right)=a_e
\tag{6.3}\label{eq:coeffE}
\]
and fixed-point expectations
\[
 E_N(a)=\frac1{|G|}\sum_{g\in G}\alpha_g^{(N)}(a),
\tag{6.4}\label{eq:fixedE}
\]
the diagrams commute exactly, and their limits are the canonical
expectations $F$ and $E$.  They are faithful.  Moreover,
\[
 E_N(x)\ge\frac1{|G|}x\quad(x\ge0),
\qquad
 E(x)\ge\frac1{|G|}x\quad(x\ge0).
\tag{6.5}\label{eq:PP}
\]
Thus the uniform Pimsner--Popa index bound survives the limit.  If each
action is saturated/outer on factors, the index is exactly $|G|$.
\end{theorem}

\begin{proof}
Equivariance gives a homomorphism
\[
 j_N:A_N\rtimes G\longrightarrow A_{N+1}\rtimes G,\qquad
 j_N(a)=\iota_N(a),\quad j_N(u_g)=u_g.
\tag{6.6}\label{eq:crossmap}
\]
It is injective.  To see this without appealing to exactness, choose a
faithful representation $\pi_{N+1}$ of $A_{N+1}$ and use the regular
covariant representation on $\ell^2(G)\otimes\Hh$.  Its restriction to
$A_N\rtimes G$ is faithful because finite groups are amenable, so full
and reduced crossed products coincide.  Hence the inductive limit on the
left of \eqref{eq:crossediso} is defined by injective maps.

Let $A_{\rm alg}=\bigcup_N\iota_{N,\infty}(A_N)$, dense in $A$.
Every algebraic crossed-product element
$x=\sum_{g\in G}a_gu_g$ with $a_g\in A_{\rm alg}$ lies at one common
finite stage because $G$ is finite.  Therefore the union of the images
of $A_N\rtimes_{\alpha^{(N)}}G$ is the algebraic crossed product
$A_{\rm alg}\rtimes G$, which is dense in $A\rtimes_\alpha G$.
The regular representations used above show that the norm of such an
$x$ is the same whether computed at its finite stage, in the inductive
limit, or in $A\rtimes_\alpha G$.  The dense isometric homomorphism
therefore extends to \eqref{eq:crossediso}.

On algebraic elements,
\[
 F_{N+1}j_N\!\left(\sum_ga_gu_g\right)
 =\iota_N(a_e)=\iota_NF_N\!\left(\sum_ga_gu_g\right),
\]
and \eqref{eq:equivariance} gives
$E_{N+1}\iota_N=\iota_NE_N$.  Hence both expectations descend to the
limit and agree there with \eqref{eq:coeffE} and
\eqref{eq:fixedE}.  Faithfulness of $F$ follows from the faithful regular
representation.  Faithfulness of $E$ follows because a sum of positive
operators can vanish only when every summand vanishes.

Finally, for $x\ge0$,
\[
 E_N(x)=\frac1{|G|}
 \left(x+\sum_{g\ne e}\alpha_g^{(N)}(x)\right)
 \ge\frac1{|G|}x.
\]
The same direct formula proves the limit inequality.  For outer finite
group actions on factors, the fixed-point index theorem gives equality
of the index with $|G|$.
\end{proof}

\subsection{State preservation and the charged grades}

\begin{corollary}[State-preserving expectations commute with the limit]
\label{cor:stateE}
Let $\varphi_N$ be compatible $\alpha^{(N)}$-invariant states.  Then
$\varphi_N\circ E_N=\varphi_N$, the limit state $\varphi$ satisfies
$\varphi\circ E=\varphi$, and the GNS extensions of $E_N$ converge on
the dense inductive union to the GNS extension of $E$.
\end{corollary}

\begin{proof}
Invariance gives
\[
 \varphi_N(E_N(a))
 =|G|^{-1}\sum_g\varphi_N(\alpha_g^{(N)}(a))
 =\varphi_N(a).
\]
Compatibility permits passage to the dense union.  Contractivity of
conditional expectations extends the equality and convergence by
continuity.
\end{proof}

\begin{corollary}[$\Z_4$ fusion and expectation]\label{cor:Z4}
For $G=\Z_4$, the limiting homogeneous subspaces
$\B_k=\overline{Au^k}$ satisfy
\[
 \B_k\B_\ell\subset\B_{k+\ell},\qquad
 \B_k^*=\B_{-k},\qquad u^4=\id.
\tag{6.7}\label{eq:grades}
\]
Thus the grade-one field obeys
$\Psi_\lambda^4\in\A$ and
$\Psi_\lambda^2\in\B_2$, identified by
\eqref{eq:lambda-order} with the vector/parity class.  The limit
expectation is
\[
 E(x)=\frac14\sum_{r=0}^3\alpha_r(x),\qquad E(x)\ge\frac14x
\quad(x\ge0).
\tag{6.8}\label{eq:quarterE}
\]
\end{corollary}

\begin{proof}
For $a,b\in A$,
$(au^k)(bu^\ell)=a\alpha_k(b)u^{k+\ell}$ and
$(au^k)^*=\alpha_{-k}(a^*)u^{-k}$.  The remaining statements follow from
Theorem~\ref{thm:crossedlimit}.
\end{proof}

\begin{remark}[Inner finite shadows versus outer type-III actions]
At a finite matrix stage the clock action can be inner, in which case
$A_N\rtimes\Z_4\cong A_N\otimes C^*(\Z_4)$ and the inclusion need not
model an irreducible outer fixed-point subfactor.  The limit theorem
above remains valid, but index \emph{equality} and Q-system
irreducibility require outerness/saturation of the limiting local action.
Those properties belong to L2/\textup{(ALG-EXH)}, not to the formal
crossed-product exchange.
\end{remark}

\section{L4: holomorphy and identification with \texorpdfstring{$(E_8)_1$}{E8 level 1}}
\label{sec:L4}

\subsection{The neutral net and its index}

The intended neutral current net is
\[
 \A=\A_{D_5,1}\otimes\A_{A_3,1}.
\tag{7.1}\label{eq:basenet}
\]
Its central charge is
\[
 c(D_5,1)+c(A_3,1)=5+3=8.
\tag{7.2}\label{eq:c8}
\]
The discriminant groups of $D_5$ and $A_3$ each have order four, so the
product lattice net has
\[
 \mu_\A=|\det D_5|\,|\det A_3|=4\cdot4=16.
\tag{7.3}\label{eq:mu16}
\]
The isotropic diagonal subgroup generated by $\lambda$ has order four.
The exact lattice identities \eqref{eq:lambda-norm}--\eqref{eq:census}
verify integer spin, evenness and unimodularity of the extension lattice.

\begin{theorem}[Identification once L1/L2 hold]\label{thm:identification}
Assume L1 and L2 produce a local irreducible index-four extension
$\A\subset\B$ with the common stress tensor, with the four grades
generated by $\lambda$, and with finite-dimensional $L_0$ eigenspaces.
Then $\B$ is the $(E_8)_1$ conformal net.
\end{theorem}

\begin{proof}
\emph{Complete rationality and holomorphy.}
The affine lattice net \eqref{eq:basenet} is completely rational, split
and strongly additive.  The order-four local simple-current Q-system has
dimension/index four by Theorem~\ref{thm:Qsystems}.  KLM therefore gives
\[
 \mu_\B=\frac{\mu_\A}{[\B:\A]^2}
 =\frac{16}{4^2}=1.
\tag{7.4}\label{eq:muone}
\]
Hence $\B$ is holomorphic.

\emph{Conformal covariance and split property.}
The common stress tensor is the Osborne--Stottmeister limit, and L2
assumes its covariance acts on the charged primary fields.  Thus the
extension is conformal, not merely M\"obius covariant.  L1 supplies the
energy-bounded generating fields.  Their finite-energy spaces reconstruct
a simple strongly local unitary VOA $V_\B$ by the CKLW/Fredenhagen--J\"or{\ss}
correspondence.  The character below is finite for every
$\operatorname{Im}\tau>0$, so \eqref{eq:traceclass} holds and
Theorem~\ref{thm:split} supplies the split property independently.

\emph{The lattice and character.}
The extension lattice
\[
 L=\bigcup_{k=0}^3(D_5\oplus A_3+k\lambda)
\tag{7.5}\label{eq:gluedlattice}
\]
is even by $\lambda\cdot(D_5\oplus A_3)\subset\Z$ and
$\norm{\lambda}^2=2$.  Its determinant is
\[
 \det L=\frac{\det(D_5)\det(A_3)}{4^2}
 =\frac{16}{16}=1.
\tag{7.6}\label{eq:unimodular}
\]
There is a unique positive-definite even unimodular lattice of rank
eight, namely $E_8$.  Therefore $L\cong E_8$.
Its theta series is the unique normalized weight-four modular form,
\[
 \Theta_L=E_4,
\tag{7.7}\label{eq:thetaE4}
\]
and
\[
 \operatorname{ch}_{V_L}
 =\frac{\Theta_L}{\eta^8}
 =\frac{E_4}{\eta^8}.
\tag{7.8}\label{eq:character}
\]
The exact initial coefficients are
$1,248,4124,34752$; the $248$ weight-one states are the $240$ roots in
\eqref{eq:census} plus eight Cartan currents.

\emph{Uniqueness.}
The reconstructed VOA is simple, unitary, rational, $C_2$-cofinite and
holomorphic: rationality and $C_2$-cofiniteness follow here from the
finite simple-current extension of the positive-definite lattice VOA.
It has $c=8$.  Theorem~\ref{thm:DM} gives
$V_\B\cong V_{E_8}$, and CKLW identifies the generated net with
$\A_{V_{E_8}}$, the level-one $E_8$ current net.
\end{proof}

\begin{remark}[Which hypotheses are presently established]
At finite level, the order-four class, integer spin, census, determinant,
character coefficients, $k=1$ current normalization, and measured
$c=8.004666$ are established.  At continuum level, the neutral
fermion-bilinear currents and stress tensor lie in the
Osborne--Stottmeister theorem class.  What is not established is that the
non-bilinear lattice charged fields satisfy L1/L2 and exhaust the local
scaling algebra.  Without that step, the finite $E_4$ character match is
not a proof of conformal covariance or splitness of the collar limit.
\end{remark}

\begin{corollary}[Local factor type]\label{cor:typeIII}
Under Theorem~\ref{thm:identification}, every nontrivial interval algebra
$\B(I)$ is the hyperfinite type-$\mathrm{III}_1$ factor.
\end{corollary}

\begin{proof}
The $E_8$ level-one current net is split, irreducible, strongly additive
and conformal.  The standard type theorem for nontrivial conformal nets
gives type $\mathrm{III}_1$; splitness and the current-net construction
give hyperfiniteness.
\end{proof}

\section{Downstream corollary cascade}

\begin{corollary}[Charged-field and $\Psi_\lambda$ convergence]
If the final gap box is discharged, the field
$\Psi_{\lambda,N}(f)$ converges on the finite-energy core, is relatively
local, has order-four fusion, and generates the missing odd
$64+64=128$ weight-one states.  The conditional statement in
\emph{$\Psi_\lambda$ at the Quasi-Free Boundary} becomes unconditional
for this collar.
\end{corollary}

\begin{proof}
Apply L1, L2 and Corollary~\ref{cor:Z4}; the state count is
\eqref{eq:census}.
\end{proof}

\begin{corollary}[Type-III KMS uniqueness]
If the final gap box is discharged, the hypotheses of the abstract
holomorphic rotational KMS theorem in
\emph{The Holomorphic Rotational KMS Extension Theorem} hold for the
concrete collar limit.  Thus the rotational $\beta$-KMS state of the
extension is the unique sector Gibbs state, and the canonical
$\Z_4$-invariant extension of the subnet state is
$\omega_\beta\circ E$.
\end{corollary}

\begin{proof}
Theorem~\ref{thm:identification} gives a holomorphic conformal net with
$\mu=1$ and Corollary~\ref{cor:typeIII} gives the type-III local factors.
The expectation is Corollary~\ref{cor:Z4}.  KLM leaves only the vacuum
DHR sector; the rotational KMS classification then leaves only its Gibbs
state.
\end{proof}

\begin{corollary}[P1 and Quillen premises: exact conditional scope]
If, in addition, the determinant-line modular response is identified
with an additive transitive four-channel response and the angular modular
normalization is $\beta_{\rm angle}=2\pi$, then
\[
 c_3=\frac1{4\beta_{\rm angle}}=\frac1{8\pi}.
\tag{8.1}\label{eq:P1}
\]
The KMS state has no independent sector-mixing direction, supplying the
rigidity premise used in the Quillen argument.
\end{corollary}

\begin{proof}
The unique invariant state assigns equal weights to four transitively
permuted channels, hence weight $1/4$ to each.  Multiplying the total
response $1/\beta_{\rm angle}$ by $1/4$ proves \eqref{eq:P1}.
KMS uniqueness removes a second state direction.
\end{proof}

\begin{remark}
The master theorem would not by itself identify the determinant-line
response with the channel response.  That is a separate
P1/Quillen operator-identification premise.  Accordingly the master
theorem closes the \emph{net/KMS rigidity premise}, not the full Quillen
variation formula.
\end{remark}

\section{Per-lemma status}

\renewcommand{\arraystretch}{1.15}
\begin{longtable}{@{}p{0.13\textwidth}p{0.24\textwidth}p{0.16\textwidth}p{0.39\textwidth}@{}}
\toprule
Item & Statement & Status & Reason\\
\midrule
\endhead
S0 & Frozen generator arithmetic:
$\norm{\lambda}^2=2$, $h_\lambda=1$, order four, census
$[52,64,60,64]$ & \mP &
Exact rational lattice calculation; rechecked 10/10.\\
S1 & Neutral bilinear currents and stress tensor have a scaling limit &
\mR &
Osborne--Stottmeister Theorems 4.11, 4.16 and 5.1; the frozen neutral
objects are in the bilinear class.\\
S2a & Continuum $e^\lambda$ has a polynomial energy bound &
\mP &
It is a norm-two, weight-one affine current; Lemma~\ref{lem:continuumEB}
and CKLW.\\
UGF & Edge-projected Green domination and local Sobolev convergence &
\mP &
Theorem~\ref{thm:UGF}: exact particle--hole Schur count plus the sharp
sine comparison gives $K_G=\pi^2/4$ for every $N$.\\
L1 & The finite fields have an $N$-uniform energy bound &
\mP &
Theorem~\ref{thm:L1reduction}, using the proved UGF and norm-one
cocycle/charge shifts.\\
TEL-a & Dyadic edge-mode isometries and charged graph defect &
\mP &
Theorem~\ref{thm:TELa}: explicit second-quantized Shannon embeddings,
$p=3$, $\epsilon=1$, and rate $N^{-2}$ for fields and adjoints.\\
TEL-b & Full-CAR two-edge remainder has an all-$N$ bound &
\mP &
v1022/v1025 inputs plus v1026 CF/DG, one-sided CROSS, validated Fourier
core and full aliases prove $\norm{R_N}_{\HS}<2.995906<3$ for all even
$N\ge16$. Norm only; no one-edge algebra identification or historical
$2.11$ theorem is claimed.\\
L2a & Edge-projected strong convergence on the finite-energy core &
\mP &
The dyadic telescoping series converges by Theorem~\ref{thm:TELa}.\\
ALG-rel & Relative locality with the neutral current net &
\mP &
Theorem~\ref{thm:relativeLocality}: integral root pairings plus strong
and graph convergence pass the exact finite commutator to the limit.\\
ALG-out & Outerness of the limiting $\Z_4$ action &
\mP &
Theorem~\ref{thm:outerness}: the connected multiplicity-free Q-system
gives $\A(I)'\cap\B^{\rm gen}(I)=\C$.\\
ALG-exh & Local algebra convergence and no extra UV operators &
\mG &
Reduced exactly to the two Buchholz--Verch estimates
\eqref{eq:EXHsharp}.\\
L3a & Crossed product commutes with the inductive limit &
\mP &
Theorem~\ref{thm:crossedlimit}, proved isometrically for every finite
group.\\
L3b & Expectations and uniform index bounds commute with the limit &
\mP &
Equations \eqref{eq:coeffE}--\eqref{eq:PP} and
Corollary~\ref{cor:stateE}.\\
L4a & Index-four extension has $\mu=1$ &
\mR &
KLM, after L2 supplies a local irreducible conformal Q-system:
$16/4^2=1$.\\
L4b & Holomorphic $c=8$ extension is $(E_8)_1$ &
\mR &
Even-unimodular rank-eight classification, Dong--Mason and CKLW,
after L1/L2 supply the strongly local VOA/net hypotheses.\\
L4c & Limit local algebras are hyperfinite type $\mathrm{III}_1$ &
\mR &
Standard conformal-current-net factor theorem after L4b.\\
KMS & Unique rotational KMS extension &
\mR &
The downstream theorem applies once L4 is available.\\
P1/Q & Equal four-channel weight and Quillen rigidity premise &
\mR &
Net/KMS part follows; determinant-line response identification remains
separate.\\
\bottomrule
\end{longtable}

\section{The remaining analytic gap}

\begin{center}
\setlength{\fboxsep}{10pt}
\fcolorbox{red!65!black}{red!3}{
\begin{minipage}{0.92\textwidth}
\textbf{Remaining microscopic-exhaustion gaps (status 2026-09-05).}

\smallskip
The Green-form and edge-mode graph parts are no longer assumptions:
Theorems~\ref{thm:UGF} and \ref{thm:TELa} prove them with
$K_G=\pi^2/4$ and dyadic rate $N^{-2}$.  Round-3 certificates close the
interval-TV input, C1 and C2a.  The analytic scalar part of C2b is
$c_{\rm potential}=0.399224424139106\ldots<0.400$.  This does not yet
identify the edge-mode system with the sharply localized scaling algebra
of the sixteen-Majorana $M=1$, $N_y=8$ microscopic collar.

\smallskip\noindent
\textbf{(TEL-B) Relaxed norm obligation closed in round 4.}
For every even $N\ge16$, the now-proved estimates are
\[
 \sqrt N\,\norm{(K'_N-K_N)_{\rm offdiag}}_{\HS}<0.065697,
 \qquad
 \sqrt N\,\norm{(\mathrm{Res}_N-K_N)_{\rm diag}}_{\HS}<0.042959.
\]
The two components are Hilbert--Schmidt orthogonal.  Combined with the
certified scalar contribution, these give $c_E<0.466902<0.542702$.
The one-sided CROSS bound is $0.282637/N$, giving Bound-B squared
$<0.714386$. A 512-node Acb certificate with infinite alias control
proves the weighted Fourier core $<1.783260551$; the actual finite-grid
smooth remainder additionally costs $<0.017333$ in aliases and is
$<1.800594$, not the former unsupported $1.7833$. With the exact jump
norm $\sqrt2$, the unrounded assembly gives $\norm{R_N}_{\HS}<2.995906<3$.
The complete all-$N$ proof is the round-4 supplement in
\texttt{tfpt\_research\_contracts.tex}, executed by
\texttt{v1026\_telb\_round4\_closure.py}. Neither the historical sharper
$2.11$ target nor the absolute-CROSS $0.283/N$ mutant is proved.
The dyadic-increment route is retired, not an alternative premise.

\smallskip\noindent
\textbf{(ALG-EXH) Buchholz--Verch upper inclusion.}
Prove the two estimates \eqref{eq:EXHsharp}; equivalently, construct a
support-preserving microscopic realization of the dyadic embeddings and
prove, for every interval $I$,
\[
 \limsup_N\B_N(I)
 \subset
 W^*\!\left(\A(I),
 \overline{\Psi_{k\lambda}(f)}:
 k\in\Z_4,\ \supp f\subset I\right)
 \subset
 \liminf_N\B_N(I).
\]
Equivalently, no bounded local ultraviolet sequence survives outside
the neutral currents and the four charged grades.  Relative locality
and outerness are already Theorems~\ref{thm:relativeLocality} and
\ref{thm:outerness}; they are not part of this residue.  The binding
ALG2 correction above means that full-CAR Vandermonde interpolation does
not prove this inclusion for $\B_N^{\rm gen}$; FE-GEN or an exact
gauge/code-projected microscopic-algebra theorem is independently required.

\smallskip
For the edge-projected theorem, \textup{(ALG-EXH)} remains the sole hypothesis.
For the original microscopic-collar formulation, the CF/DG norm inputs are
now proved by the Round-4 certificate (v1026), for the fixed-width
$M=1$, $N_y=8$ model and every even $N\ge16$.
This does not supply the microscopic-algebra identification required by
FE-GEN or prove \textup{(ALG-EXH)}.
Once \eqref{eq:EXHsharp} holds, L1/L2 are complete and L3 plus the cited
L4 chain finish the theorem.
\end{minipage}}
\end{center}

\paragraph{Round-7 microscopic twist boundary (2026-09-05).}
The repository checker \texttt{v1033\_charged\_disorder.py} constructs
order-four charged corners in the actual QWZ matter Fock space only after
adding a local four-level flux register to the seam Hamiltonian.
An onsite scalar matter gauge transformation alone cannot change its
Wilson holonomy. The eight-channel quarter twist has charge
$q=(1/4)^8$, hence $h=\lVert q\rVert^2/2=1/4$, not the required
$h_\lambda=1$. Every same-phase integer dressing with norm squared two
has odd fermion parity. The $D_5$ half-twist and $A_3$ inverse-quarter
description additionally requires its diagonal-$U(1)$ quotient and
cocycle/GSO identification; it is not automatically the same microscopic
eight-channel collar. The source's forward seam $i^r$ fixes the Bloch
shift to $k=(2\pi n-r\pi/2)/N$ in the stated convention.
Finite sharp-string energy and nonlocal smooth-twist scaling are
diagnostics, not a conformal-weight or support-preserving FE-GEN proof.
Thus the construction supplies actual corners in an enlarged model,
not the missing MMST upper inclusion. Complete definitions, proofs and
typed regressions are in the round-7 supplement of
\texttt{tfpt\_research\_contracts}.

\section*{Conclusion}
\addcontentsline{toc}{section}{Conclusion}

The order-four arithmetic and the finite-group operator algebra are not
the obstruction.  The edge-projected Green form is now uniformly
dominated for every $N$ with $K_G=\pi^2/4$, and it converges to the
chiral Green form in local $H^2$ quadratic-form norm.  Exact compatible
dyadic edge-mode isometries now give the charged graph-defect rate
$N^{-2}$ for fields and adjoints.  The full
crossed-product/expectation limit is proved, and the holomorphic $c=8$
endpoint is rigid once a local conformal limit with the charged field
exists.  The unresolved point is the support-preserving identification
of this edge-mode limit with the full microscopic CAR local scaling
algebra. For the original microscopic collar, v1026 now proves the
relaxed uniform TEL-b norm below $2.995906<3$, including CF, DG and
finite-grid aliases, without a finite-plateau inference. Relative locality and
outerness are proved.  The remaining in-house algebraic obligation is
the Buchholz--Verch upper inclusion with the binding FE-GEN/gauge-code
qualification. The conditional sector/actual-word frame theorem in v1030
is a concrete sufficient interface, not a proof that the microscopic
one-edge algebra satisfies it. Joint-adjoint control and high-output tails
remain indispensable.

\begin{thebibliography}{99}

\bibitem{BDL}
D.~Buchholz, C.~D'Antoni and R.~Longo,
\emph{Nuclear maps and modular structures II: applications to quantum
field theory},
Commun.\ Math.\ Phys.\ \textbf{129} (1990), 115--138.

\bibitem{BuchholzWichmann}
D.~Buchholz and E.~H.~Wichmann,
\emph{Causal independence and the energy-level density of states in
local quantum field theory},
Commun.\ Math.\ Phys.\ \textbf{106} (1986), 321--344.

\bibitem{BuchholzVerch}
D.~Buchholz and R.~Verch,
\emph{Scaling algebras and renormalization group in algebraic quantum
field theory},
Rev.\ Math.\ Phys.\ \textbf{7} (1995), 1195--1239.

\bibitem{CKLW}
S.~Carpi, Y.~Kawahigashi, R.~Longo and M.~Weiner,
\emph{From vertex operator algebras to conformal nets and back},
Mem.\ Amer.\ Math.\ Soc.\ \textbf{254} (2018), no.~1213,
\href{https://arxiv.org/abs/1503.01260}{arXiv:1503.01260}.

\bibitem{CarpiTomassini}
S.~Carpi and T.~Tomassini,
\emph{Energy bounds for vertex operator algebra extensions},
Lett.\ Math.\ Phys.\ \textbf{113} (2023), 87,
\href{https://arxiv.org/abs/2303.14097}{arXiv:2303.14097}.

\bibitem{DongMason}
C.~Dong and G.~Mason,
\emph{Holomorphic vertex operator algebras of small central charge},
Pacific J.\ Math.\ \textbf{213} (2004), 253--266,
\href{https://arxiv.org/abs/math/0203005}{arXiv:math/0203005}.

\bibitem{Gui}
B.~Gui,
\emph{Energy bounds condition for intertwining operators of types
$B$, $C$, and $G_2$ unitary affine vertex operator algebras},
Trans.\ Amer.\ Math.\ Soc.\ \textbf{372} (2019), 7371--7424,
\href{https://arxiv.org/abs/1809.07003}{arXiv:1809.07003}.

\bibitem{KLM}
Y.~Kawahigashi, R.~Longo and M.~M\"uger,
\emph{Multi-interval subfactors and modularity of representations in
conformal field theory},
Commun.\ Math.\ Phys.\ \textbf{219} (2001), 631--669.

\bibitem{LongoRehren}
R.~Longo and K.-H.~Rehren,
\emph{Nets of subfactors},
Rev.\ Math.\ Phys.\ \textbf{7} (1995), 567--597.

\bibitem{OsborneStottmeister}
T.~J.~Osborne and A.~Stottmeister,
\emph{Conformal field theory from lattice fermions},
Commun.\ Math.\ Phys.\ \textbf{398} (2023), 219--283,
\href{https://arxiv.org/abs/2107.13834}{arXiv:2107.13834}.

\bibitem{Peller}
V.~V.~Peller,
\emph{Hankel operators of class $\mathfrak S_p$ and their applications
(rational approximation, Gaussian processes, the problem of majorizing
operators)},
Math.\ USSR Sbornik \textbf{41} (1982), 443--479.

\bibitem{ToledanoLaredo}
V.~Toledano Laredo,
\emph{Fusion of positive energy representations of $L\operatorname{Spin}_{2n}$},
Ph.D.\ thesis, University of Cambridge, 1997; analytic energy-bound
results as used in the type-$D$ loop-group literature.

\bibitem{Wassermann}
A.~Wassermann,
\emph{Operator algebras and conformal field theory III: fusion of
positive energy representations of $LSU(N)$ using bounded operators},
Invent.\ Math.\ \textbf{133} (1998), 467--538.

\end{thebibliography}

\end{document}
